% ~2 min
\titleframe{Calculating Correct and Efficient Programs}{%
  \authorblock{David Castro-Perez}{David.CastroPerez@rhul.ac.uk}
    {RHULLogo.png}{Royal Holloway, University of London}}{%
  \par\medskip
  {\footnotesize Based on joint work with Marco Paviotti and Michael Vollmer
    (University of Kent):\\ \emph{Program Optimisations via Hylomorphisms
    for Extraction of Executable Code}, ITP 2025.}%
  \speakernote{Joint work with Marco Paviotti and Michael Vollmer (ITP
    2025). Audience: computer scientists outside programming languages. The
    talk is about getting programs that are fast and provably correct at
    the same time, without proving each fast program from scratch.}}

\begin{frame}
  \frametitle{In One Slide}
  \begin{center}
    \begin{tikzpicture}[alt={A clear but slow specification is turned into
        a fast program by a chain of rewrite steps. Each step is a proved
        law, so the fast program provably computes the same results as the
        specification.},
      stage/.style = {text width = 30mm, align = center,
        minimum height = 14mm}]
      \node[box, stage] (s) {\textbf{clear, slow}\\ specification};
      \node[hot, stage, right = 50mm of s] (p) {\textbf{fast}\\ program};
      \draw[->, very thick] (s) -- node[above, note, align = center]
        {rewrite with\\ \emph{proved} laws} (p);
    \end{tikzpicture}
  \end{center}
  \begin{columns}
    \begin{column}{0.45\textwidth}
      \begin{sticky}
        \textbf{The idea} (1980s):
        \begin{itemize}
          \item Programs as algebra
          \item Optimise by \underline{rewriting}
          \item Correct by construction
        \end{itemize}
      \end{sticky}
    \end{column}
    \hspace{5mm}
    \begin{column}{0.45\textwidth}
      \begin{sticky}
        \textbf{Our work:}
        \begin{itemize}
          \item Every step \underline{machine-checked}
          \item No extra axioms
          \item Runnable code out
        \end{itemize}
      \end{sticky}
    \end{column}
  \end{columns}
  \speakernote{The whole talk in one picture. Start from a specification
    that is obviously right, even if slow. Rewrite it, step by step, with
    algebraic laws; each step is an equation, so the end result equals the
    specification. This is an old idea (Bird, Meertens: the Algebra of
    Programming). What we add: every step is checked by a proof assistant
    (Rocq), with no extra assumptions, and the result is a program you can
    run. At the end: where we want to take this, down to low-level code.}
\end{frame}
